% How to Survive Your First Corporate Job - Course Slides
% Compile with: pdf, then pdflatex (for references)
\documentclass[aspectratio=169, 11pt]{beamer}
\usetheme{Madrid}
\usecolortheme{whale}
\usefonttheme{professionalfonts}

% Required packages
\usepackage[utf8]{inputenc}
\usepackage{amsmath, amssymb}
\usepackage{graphicx}
\usepackage{booktabs}
\usepackage{hyperref}
\usepackage{xcolor}

% Custom colors (corporate but modern)
\definecolor{primary}{RGB}{0, 82, 112}    % deep teal
\definecolor{accent}{RGB}{200, 83, 0}     % burnt orange
\definecolor{softgray}{RGB}{245, 247, 250}
\setbeamercolor{title}{fg=white, bg=primary}
\setbeamercolor{frametitle}{fg=primary, bg=softgray}
\setbeamercolor{block title}{fg=white, bg=primary}
\setbeamercolor{block body}{fg=darkgray, bg=softgray}

% Remove navigation symbols
\setbeamertemplate{navigation symbols}{}
\setbeamertemplate{footline}[frame number]

% Title slide styling
\title{\textbf{How to Survive Your in The Uncertain Remote Job}}
\subtitle{Learning When Nobody Teaches You}
\author{Professional Development}
\date{}

\begin{document}
	
	% ----------------------------------------------------------
	% SLIDE 1: TITLE
	% ----------------------------------------------------------
	\begin{frame}
		\titlepage
		\centering
		\vspace{0.5cm}
		\small{Navigate uncertainty · Learn independently · Thrive without hand-holding}
	\end{frame}
	
	
	
	% ----------------------------------------------------------
	% SLIDE 1: THE UNIVERSITY VS. REALITY GAP (Finance/Quant Focus)
	% ----------------------------------------------------------
	\begin{frame}{The gap: university vs. the trading floor / risk desk}
		\begin{block}{What university taught you to expect}
			\begin{itemize}
				\item Clean, complete datasets
				\item Step‑by‑step documentation
				\item Clear problem statements
				\item One correct answer
				\item A manager who explains your task
			\end{itemize}
		\end{block}
		\vspace{0.2cm}
		\alert{\textbf{What the job actually gives you:}}
		\begin{itemize}
			\item Incomplete or missing information — \textit{day one and beyond}
			\item Systems with \textbf{little to no documentation} (you search yourself)
			\item Confusing processes and \textbf{unclear ownership}
			\item Multiple teams, deadlines, constant changes — \textit{parallel tasks}
		\end{itemize}
		\vspace{0.1cm}
		\textit{No one warns you. This course closes the gap.}
	\end{frame}
	
	% ----------------------------------------------------------
	% SLIDE 2: WHAT THIS COURSE TEACHES (BEYOND TECHNICAL MODELS)
	% ----------------------------------------------------------
	\begin{frame}{What this course focuses on — the unwritten skills}
		\begin{center}
			\large{Not just another modeling class.}
		\end{center}
		\vspace{0.2cm}
		\textbf{You learn how to:}
		\begin{columns}[T]
			\column{0.5\textwidth}
			\begin{itemize}
				\item Navigate chaos in documentation \& systems
				\item Work effectively with uncertainty \& ambiguity
				\item Understand your role when guidance is limited
			\end{itemize}
			\column{0.5\textwidth}
			\begin{itemize}
				\item Handle cross‑team confusion \& communication gaps
				\item Build confidence without full understanding
				\item Operate in a \textbf{managerless}, self‑directed way
			\end{itemize}
		\end{columns}
		\vspace{0.3cm}
		\alert{For roles in:}
		\begin{itemize}
			\item Financial analytics
			\item Risk / AML / model validation
			\item Data / quant roles in large organizations
		\end{itemize}
		\vspace{0.1cm}
		\textit{The skills that separate juniors from experienced pros.}
	\end{frame}
	
	% ----------------------------------------------------------
	% SLIDE 3: HOW YOU ACTUALLY SURVIVE (TACTICAL, NO FLUFF)
	% ----------------------------------------------------------
	\begin{frame}{The survival workflow — when no one teaches you}
		\begin{enumerate}
			\item \textbf{Find the real source of truth} \\
			\small(outdated wiki? Slack history? that one Excel file? the person who knows)
			
			\item \textbf{Work with incomplete information} \\
			\small(identify what would break everything → test small bets → reversible decisions)
			
			\item \textbf{Handle unclear ownership} \\
			\small(who benefits most? volunteer temporarily + propose permanent owner)
			
			\item \textbf{Ask better questions} \\
			\small(not “how do I do this?” — instead: “I tried X and Y, stuck on step 3, can you point me to the database?”)
			
			\item \textbf{Document as you go} \\
			\small(personal knowledge base: OneNote/Obsidian + timestamped findings)
		\end{enumerate}
		\vspace{0.2cm}
		\alert{Golden rule for banks \& large finance:} \\
		Never escalate a problem without suggesting at least one possible fix.
	\end{frame}
	
	
	
	% ==========================================================
	% SECTION 2: DOCUMENTATION HELL (Shadow Analyst + Independent Working Focus)
	% ==========================================================
	\section{Documentation Hell}
	
	% ----------------------------------------------------------
	% SLIDE 2.1: READING BAD DOCUMENTATION (Independent Worker Edition)
	% ----------------------------------------------------------
	\begin{frame}{2.1 Reading bad documentation — no one hands you a manual}
		\begin{block}{What you actually find (remote + managerless reality)}
			\begin{itemize}
				\item Documents with \textbf{missing logic} — steps just vanish
				\item \textbf{Outdated information} — last updated before you joined
				\item \textbf{Inconsistent instructions} — page 3 says one thing, page 10 another
				\item No owner, no version, no date stamp
				\item \textit{And no manager to ask — you figure it out alone}
			\end{itemize}
		\end{block}
		\vspace{0.1cm}
		\alert{\textbf{Shadow analyst / independent working skill:}}
		\begin{enumerate}
			\item Find one working example → trace backward to understand logic
			\item Compare with old emails, Slack history, or shared drive versions
			\item Infer missing steps from context — write down your assumption
			\item Convert your findings into \textbf{wiki-style notes} for the whole team
		\end{enumerate}
		\textit{You become the person who fixed the documentation, not just read it.}
	\end{frame}
	
	% ----------------------------------------------------------
	% SLIDE 2.2: WRITING DOCUMENTATION FOR BUSY LEADERS (C-Suite Ready)
	% ----------------------------------------------------------
	\begin{frame}{2.2 Writing documentation — what executives actually want}
		\begin{columns}[T]
			\column{0.48\textwidth}
			\textbf{What most people write:}
			\begin{itemize}
				\item 20 pages of dense text
				\item Buried conclusions
				\item No executive summary
				\item Jargon everywhere
			\end{itemize}
			\column{0.48\textwidth}
			\textbf{What C-suite \& managers want:}
			\begin{itemize}
				\item \textbf{One-page summary first}
				\item Status: red / yellow / green
				\item \textit{What do I need to act on?}
				\item Clear next steps
			\end{itemize}
		\end{columns}
		\vspace{0.2cm}
		\textbf{Your shadow analyst document structure:}
		\begin{enumerate}
			\item \textbf{Executive summary} (1 page — for busy leaders)
			\item What this document covers (1 paragraph)
			\item Key findings or decisions needed
			\item Supporting details (appendix for deep divers)
		\end{enumerate}
		\alert{If a leader can't understand it in 60 seconds, rewrite it.} \\
		\textit{Your job as a shadow analyst: make busy people smarter in 5 minutes.}
	\end{frame}
	
	% ----------------------------------------------------------
	% SLIDE 2.3: SMALL CHANGE = BIG PAPERWORK (Real Remote Work Reality)
	% ----------------------------------------------------------
	\begin{frame}{2.3 Small change, big paperwork — real example}
		\begin{center}
			\Large{\textbf{One small update → hours of documentation}}
		\end{center}
		\vspace{0.2cm}
		\begin{block}{Real scenario you will face (remote, no training)}
			You're asked (via a ticket or email) to update a process or report: \\
			\textit{``Change the cutoff date from Dec 31 to Jan 31.''}
		\end{block}
		\vspace{0.1cm}
		\textbf{What actually happens — and you must track it yourself:}
		\begin{itemize}
			\item Make the actual change (1 minute)
			\item Update the process documentation (1 hour)
			\item Create a change log entry (15 min)
			\item Notify everyone who uses it (emails, Teams, Trello)
			\item Update any training or reference guides (1 hour)
		\end{itemize}
		\alert{This is normal in any large org. No one reminds you. You manage this independently.}
	\end{frame}
	
	% ----------------------------------------------------------
	% SLIDE 2.4: VERSION CONTROL & CHANGE LOGS (No One Teaches This)
	% ----------------------------------------------------------
	\begin{frame}{2.4 Version control — survive without a manager}
		\begin{block}{What amateurs do (you'll inherit this mess)}
			\begin{itemize}
				\item \texttt{report\_FINAL.docx}
				\item \texttt{report\_FINAL\_v2.docx}
				\item \texttt{report\_FINAL\_REALFINAL.docx}
			\end{itemize}
		\end{block}
		\vspace{0.1cm}
		\begin{block}{What independent pros do (audit-proof, no one tells you)}
			\begin{itemize}
				\item Use SharePoint, Google Drive, or Git with version history
				\item Change log table in \textbf{every important document}:
				\small{\begin{tabular}{lll}
						\toprule
						\textbf{Date} & \textbf{Change} & \textbf{Why} \\
						\midrule
						Jan 10 & Updated cutoff date & New fiscal calendar \\
						\bottomrule
				\end{tabular}}
				\item Each change = new version number (v1.2.0)
			\end{itemize}
		\end{block}
		\alert{No one will ask you to do this. You do it anyway. That's independent working.}
	\end{frame}
	
	% ----------------------------------------------------------
	% SLIDE 2.5: UNDERSTANDING HOW WORK FLOWS (Most Juniors Don't Know This)
	% ----------------------------------------------------------
	\begin{frame}{2.5 How work actually flows — the lifecycle nobody explains}
		\begin{center}
			\Large{Task Created → You Work → Review → Approved → Archived}
		\end{center}
		\vspace{0.3cm}
		\begin{columns}[T]
			\column{0.23\textwidth}
			\textbf{1. Task} \\
			\small{Ticket, email, Slack}
			\column{0.23\textwidth}
			\textbf{2. Do} \\
			\small{You figure it out}
			\column{0.23\textwidth}
			\textbf{3. Review} \\
			\small{Manager or peer}
			\column{0.23\textwidth}
			\textbf{4. Archive} \\
			\small{Docs + notes saved}
		\end{columns}
		\vspace{0.3cm}
		\alert{Most people focus only on step 2.} \\
		\textit{As a shadow analyst, you live in steps 1, 3, and 4 — tracking, summarizing, documenting.}
		
		\vspace{0.2cm}
		\small{If you don't understand the full cycle, you'll lose context, miss approvals, and fail to archive. Then next time, you start from zero.}
	\end{frame}
	
	% ----------------------------------------------------------
	% SLIDE 2.6: SUMMARY — DOCUMENTATION HELL SURVIVAL CHEAT SHEET
	% ----------------------------------------------------------
	\begin{frame}{Documentation Hell — the independent worker cheat sheet}
		\begin{center}
			\small
			\begin{tabular}{p{4cm} p{7cm}}
				\toprule
				\textbf{Problem} & \textbf{Survival Tactic} \\
				\midrule
				Outdated / missing docs & Reverse-engineer from working examples + timestamp everything \\
				Inconsistent instructions & Flag them clearly in your notes — don't ignore \\
				Small change = big paperwork & Maintain a change log before they ask \\
				No version control & Use SharePoint / Drive + version numbers \\
				Don't know the workflow & Map your task to: Task → Do → Review → Archive \\
				Busy leaders won't read long docs & Lead with a 1-page executive summary \\
				\bottomrule
			\end{tabular}
		\end{center}
		\vspace{0.3cm}
		\alert{The person who documents well becomes indispensable.} \\
		\textit{Everyone else gets the same questions every single week.}
	\end{frame}
	%
	
	
	% ----------------------------------------------------------
	% SLIDE: DEBUGGING LATEX (AND ANY CODE) WITH LIMITED INFO
	% ----------------------------------------------------------
	\begin{frame}{Debugging when you can't ``just learn the code''}
		\begin{block}{The real enterprise reality}
			\begin{itemize}
				\item You don't have time to fully understand everything
				\item Documentation is missing or outdated
				\item The original author left 2 years ago
				\item Proxy errors, missing packages, Git conflicts
				\item \textit{And no one is coming to help}
			\end{itemize}
		\end{block}
		\vspace{0.2cm}
		\alert{\textbf{The pragmatic approach:}}
		\begin{enumerate}
			\item \textbf{Use common sense first} — what should the output look like?
			\item \textbf{Find working past code} — copy, paste, modify, test
			\item \textbf{Change one thing at a time} — isolate the problem
			\item \textbf{Google the error} — someone else already fixed it
			\item \textbf{Make it work} → then figure out \textit{why} later
		\end{enumerate}
		\vspace{0.1cm}
		\textit{Perfect is the enemy of done. In enterprise, ``works'' beats ``fully understood''.}
	\end{frame} ----------------------------------------------------------
	% SLIDE 2: THE REALITY CHECK
	% ----------------------------------------------------------
	\begin{frame}{The uncomfortable truth}
		\begin{block}{Most training assumes perfect conditions}
			\begin{itemize}
				\item Perfect documentation
				\item Experienced mentors available 24/7
				\item Organized processes
				\item Managers with unlimited time
			\end{itemize}
		\end{block}
		\vspace{0.3cm}
		\alert{\textbf{Reality:}} \textit{Confusing systems, outdated wikis, conflicting info, and projects you barely understand.}
	\end{frame}
	
	% ----------------------------------------------------------
	% SLIDE 3: THIS COURSE IS DIFFERENT
	% ----------------------------------------------------------
	\begin{frame}{What this course focuses on}
		\begin{center}
			\large{Skills organizations rarely teach \quad → \quad but determine success}
		\end{center}
		\vspace{0.5cm}
		\begin{itemize}
			\item Operating effectively with \textbf{incomplete information}
			\item Making progress when \textbf{guidance is limited}
			\item Keeping sanity while expectations stay \textbf{high}
		\end{itemize}
		\vspace{0.3cm}
		\alert{No fluff. No fake theory. Just workplace survival mechanics.}
	\end{frame}
	
	% ----------------------------------------------------------
	% SLIDE 4: WHO THIS IS FOR
	% ----------------------------------------------------------
	\begin{frame}{Target audience}
		\begin{columns}[T]
			\column{0.5\textwidth}
			\textbf{Perfect for:}
			\begin{itemize}
				\item Recent graduates
				\item Career switchers
				\item New joiners in large corps
				\item Analysts, consultants, devs
				\item Anyone struggling with ambiguity
			\end{itemize}
			\column{0.5\textwidth}
			\textbf{Industries:}
			\begin{itemize}
				\item Finance / Banking
				\item Technology
				\item Consulting
				\item Healthcare
				\item Government  enterprises
			\end{itemize}
		\end{columns}
	\end{frame}
	
	% ----------------------------------------------------------
	% SLIDE 5: WHAT YOU WILL LEARN (Overview)
	% ----------------------------------------------------------
	\begin{frame}{Core outcomes}
		\begin{enumerate}
			\item Navigate the first 90 days without panic
			\item Overcome imposter syndrome
			\item Learn independently (no formal training)
			\item Extract gold from bad documentation
			\item Tame chaotic systems \& tools
			\item Work across enemy — sorry, multiple teams
			\item Ask better questions, get useful answers
			\item Escalate like a pro
			\item Handle politics without losing soul
			\item Deliver results when you understand 60\% of the task
		\end{enumerate}
	\end{frame}
	
	%
	
	
	 ----------------------------------------------------------
	% SECTION 1: THE SHOCK — FIRST 90 DAYS
	% ----------------------------------------------------------
	\section{The Shock — First 20 days 60 days 90 Days}
	\begin{frame}{Section 1: The Shock — Your first 90 days}
		\begin{center}
			\Large{\textbf{Why confusion is normal}}
		\end{center}
		\vspace{0.3cm}
		\begin{itemize}
			\item Job descriptions rarely match reality
			\item Everyone feels imposter syndrome (yes, even your boss)
			\item You'll be productive while understanding only 40\% of the jargon
		\end{itemize}
	\end{frame}
	
	% ----------------------------------------------------------
	% SLIDE 7: IMPOSTER SYNDROME IS A LIAR
	% ----------------------------------------------------------
	\begin{frame}{Imposter syndrome — the silent career killer}
		\begin{block}{Typical thoughts}
			\ I don't belong here   They'll find out I'm faking it  Everyone knows more than me  
		\end{block}
		\vspace{0.2cm}
		\textbf{Reality check:}
		\begin{itemize}
			\item Most colleagues are also figuring things out
			\item Confidence comes from action, not knowledge
			\item Strategy: track small wins daily
		\end{itemize}
	\end{frame}
	
	% ----------------------------------------------------------
	% SLIDE 8: SURVIVING THE FIRST MONTH
	% ----------------------------------------------------------
	\begin{frame}{First 30 days — tactical survival}
		\begin{enumerate}
			\item Find the \textbf{one person} who knows where bodies are buried
			\item Learn the \textbf{actual} workflow (ignore org chart)
			\item Create a personal glossary of acronyms
			\item Identify top 3 recurring problems on your team
			\item Ship one tiny improvement (builds reputation)
		\end{enumerate}
		\alert{Productivity   knowing everything. Productivity = making forward progress.}
	\end{frame}
	
	% ----------------------------------------------------------
	% SECTION 2: WORKING WITHOUT GUIDANCE
	% ----------------------------------------------------------
	\section{Working Without Guidance}
	\begin{frame}{Section 2: Working without guidance}
		\begin{center}
			\large{No one has time to teach you. Now what?}
		\end{center}
		\vspace{0.4cm}
		\textbf{Self-directed learning toolkit:}
		\begin{itemize}
			\item Reverse engineer existing work products
			\item Use the   15-minute rule  (try yourself, then ask)
			\item Build a \textbf{personal knowledge base} (OneNote/Obsidian)
			\item Find internal wikis, old decks, Slack history
		\end{itemize}
	\end{frame}
	
	% ----------------------------------------------------------
	% SLIDE 10: THE ART OF ASKING (SMART QUESTIONS)
	% ----------------------------------------------------------
	\begin{frame}{Asking better questions}
		\begin{block}{Bad question}
		  How do I do this report? 
		\end{block}
		\begin{block}{Good question (shows effort)}
			 I looked at last month's template and tried X. I'm stuck on step 3 where the data source is unclear. Could you point me to the right database?  
		\end{block}
		\textbf{Why it works:} Respects others' time + proves you're not helpless.
 	\end{frame}
	
	% ----------------------------------------------------------
	% SLIDE 11: ESCALATION WITHOUT BURNING BRIDGES
	% ----------------------------------------------------------
	\begin{frame}{Escalate like a diplomat}
		\textbf{Escalation matrix:}
		\begin{enumerate}
			\item Try to unblock yourself (document your attempts)
			\item Ask immediate teammate (short, specific)
			\item Ask team lead (show what you tried)
			\item Escalate to manager with proposed solution path
		\end{enumerate}
		\vspace{0.2cm}
		\alert{Golden rule:} Never escalate a problem without suggesting at least one possible fix.
	\end{frame}
	
	% ----------------------------------------------------------
	% SECTION 3: INFORMATION CHAOS
	% ----------------------------------------------------------
	\section{Information Chaos}
	\begin{frame}{Section 3: Information chaos}
		\begin{center}
			\textit{\Large Outdated docs, 3 versions of the truth, and conflicting Slack messages}
		\end{center}
		\vspace{0.4cm}
		\textbf{Strategies:}
		\begin{itemize}
			\item Identify \textbf{source of record} (one system to rule)
			\item Timestamp your findings (document date matters)
			\item Cross-validate with 2 different people/sources
			\item Create your own \textit{living documentation}
		\end{itemize}
	\end{frame}
	
	% ----------------------------------------------------------
	% SLIDE 13: DEALING WITH CONFLICTING INFO
	% ----------------------------------------------------------
	\begin{frame}{Who do you believe? 5-step verification}
		\begin{enumerate}
			\item Check \textbf{recency} (newer tends to win, but not always)
			\item Check \textbf{authority} (process owner vs intern's note)
			\item Check \textbf{results} (what actually works in production)
			\item Ask the team in a neutral way:  I saw two approaches… 
			\item \textbf{Document the decision} for next person
		\end{enumerate}
	\end{frame}
	
	% ----------------------------------------------------------
	% SECTION 4: SYSTEMS AND TOOLS MESS
	% ----------------------------------------------------------
	\section{Systems \& Tools Mess}
	\begin{frame}{Section 4: Systems and tools mess}
		\textbf{Typical zoo:} Jira, Confluence, SharePoint, Salesforce, internal CRMs, model repositories, Teams, Slack, email threads, and that one Excel file that runs the entire department.
		\vspace{0.3cm}
		\begin{block}{Survival strategy}
			\begin{itemize}
				\item Create a \textbf{tool map} (what lives where)
				\item Learn search syntax for each platform
				\item Bookmark the top 5 most-used resources
				\item Ask:   What's the ONE tool I must master first? 
			\end{itemize}
		\end{block}
	\end{frame}
	
	% ----------------------------------------------------------
	% SLIDE 15: SHADOW THE SYSTEMS GURU
	% ----------------------------------------------------------
	\begin{frame}{Find the systems whisperer}
		Every team has one person who \textit{actually} understands how the tools connect.
		\vspace{0.3cm}
		\begin{itemize}
			\item Buy them coffee / virtual coffee
			\item Ask for a 15-min tour of \textbf{their} workflow
			\item Take notes, then share them back (they'll love you)
			\item Pay it forward — teach another new joiner
		\end{itemize}
	\end{frame}
	
	% ----------------------------------------------------------
	% SECTION 5: CROSS-TEAM CHAOS
	% ----------------------------------------------------------
	\section{Cross-Team Chaos}
	\begin{frame}{Section 5: Cross-team chaos}
		\centering
		\Large{Multiple stakeholders, different priorities, unclear ownership}
		\vspace{0.4cm}
		
		\textbf{Key skills:}
		\begin{itemize}
			\item Stakeholder mapping (who cares, who decides, who blocks)
			\item Communication style per team (legal wants detail, sales wants headlines)
			\item RACI for the clueless (Responsible, Accountable, Consulted, Informed)
		\end{itemize}
	\end{frame}
	
	% ----------------------------------------------------------
	% SLIDE 17: THE POLITICS LIGHT — BUILDING RELATIONSHIPS
	% ----------------------------------------------------------
	\begin{frame}{Organizational politics (the non-evil version)}
		Politics isn't inherently bad. It's about \textbf{influence and alignment}.
		\begin{enumerate}
			\item Be helpful before you need help
			\item Learn what each stakeholder \textit{actually} cares about (metrics, risks, visibility)
			\item Close loops: always summarize action items
			\item Never complain publicly — take it to the right person
		\end{enumerate}
	\end{frame}
	
	% ----------------------------------------------------------
	% SLIDE 18: WHEN OWNERSHIP IS UNCLEAR
	% ----------------------------------------------------------
	\begin{frame}{  That's not my job  — the enemy of progress}
		\textbf{What to do when nobody owns a task:}
		\begin{itemize}
			\item Ask: who benefits most from this being done?
			\item Volunteer \textit{temporarily} (  I can start, but need handoff after X )
			\item Create a lightweight process doc as you go
			\item After 2 weeks, propose a permanent owner
		\end{itemize}
		\alert{The person who solves orphaned work becomes invaluable.}
	\end{frame}
	
	% ----------------------------------------------------------
	% SECTION 6: DOING WORK YOU DON'T FULLY UNDERSTAND
	% ----------------------------------------------------------
	\section{Doing Work You Don't Fully Understand}
	\begin{frame}{Section 6: Doing work you do not fully understand}
		\begin{center}
			\Large{The superpower of modern professionals}
		\end{center}
		\vspace{0.3cm}
		Most systems were built years before you arrived. You'll never know 100\%.
		\vspace{0.3cm}
		\textbf{Goal:} Become effective while continuously learning.
	\end{frame}
	
	% ----------------------------------------------------------
	% SLIDE 20: DECISION-MAKING UNDER UNCERTAINTY
	% ----------------------------------------------------------
	\begin{frame}{How experienced pros decide (with 60\% clarity)}
		\begin{columns}
			\column{0.5\textwidth}
			\textbf{They don't wait for certainty:}
			\begin{itemize}
				\item Identify the \textit{one thing} that would break everything
				\item Make reversible decisions quickly
				\item Test assumptions with small bets
				\item Communicate \textbf{confidence level} (e.g., 70\% sure)
			\end{itemize}
			\column{0.5\textwidth}
			\textbf{Example:}
		 Based on current data, I recommend Option A. If assumption X is wrong, we'll fall back to B. I'll monitor daily. 
		\end{columns}
	\end{frame}
	
	% ----------------------------------------------------------
	% SLIDE 21: RISK-AWARE PRODUCTIVITY
	% ----------------------------------------------------------
	\begin{frame}{Identify risks before they bite you}
		\textbf{Simple pre-flight checklist for unfamiliar tasks:}
		\begin{enumerate}
			\item What could go wrong? (top 3 risks)
			\item Who should review this?
			\item What's the worst-case impact?
			\item How will I verify the output is correct?
			\item Do I have a rollback / undo plan?
		\end{enumerate}
		\vspace{0.2cm}
		\textit{This separates juniors from pros.}
	\end{frame}
	
	% ----------------------------------------------------------
	% SLIDE 22: TESTING ASSUMPTIONS — LIGHTWEIGHT VALIDATION
	% ----------------------------------------------------------
	\begin{frame}{The   smoke test  for corporate work}
		Before spending 10 hours on a report/model/presentation:
		\begin{itemize}
			\item Show a rough skeleton to a trusted peer (5 min feedback)
			\item Run a small subset of data through your process
			\item Ask:   Does the output pass the sanity check? 
			\item Compare with historical example (even if outdated)
		\end{itemize}
		\alert{Fail fast, fix faster, and never build in a vacuum.}
	\end{frame}
	
	% ----------------------------------------------------------
	% SLIDE 23: COMMUNICATING INCOMPLETE KNOWLEDGE
	% ----------------------------------------------------------
	\begin{frame}{How to say   I don't know  — professionally}
		\begin{block}{Weak}
			  I have no idea.  
		\end{block}
		\begin{block}{Strong (builds trust)}
			 I don't have the full picture yet. Based on what I know, my best guess is X. I'll confirm with Y and get back to you by EOD.  
		\end{block}
		\textbf{Formula:} Acknowledge gap + provide best current answer + commit to closing the gap + deadline.
	\end{frame}
	
	% ----------------------------------------------------------
	% SLIDE 24: REAL WORLD EXAMPLE — SURVIVING THE FIRST PROJECT
	% ----------------------------------------------------------
	\begin{frame}{Case study: Legacy dashboard migration}
		\textbf{Scenario:} You're asked to update a critical dashboard. Documentation is 3 years old. Original author left the company.
		\vspace{0.3cm}
		\textbf{Survival approach:}
		\begin{enumerate}
			\item Find where the data sources are actually pointing (run a test query)
			\item Compare current output vs last known correct version
			\item Document each assumption as you go
			\item Ask users:   What's one number you always check?  (smoke test)
			\item Deliver with caveats + validation plan
		\end{enumerate}
	\end{frame}
	
	
	% ==========================================================
	% SECTION: DEBUGGING WHEN NOTHING WORKS (LaTeX + General)
	% ==========================================================
	\section{Debugging When Nothing Works}
	
	% ----------------------------------------------------------
	% SLIDE 1: THE REALITY — YOU KNOW NOTHING AND NOTHING WORKS
	% ----------------------------------------------------------
	\begin{frame}{The worst-case scenario (and it will happen)}
		\begin{block}{Your situation right now:}
			\begin{itemize}
				\item No documentation — none exists
				\item No notes — previous person left nothing
				\item No working code — everything is broken
				\item No error messages you understand — just gibberish
				\item Software won't install — proxy blocks it
				\item Compilation fails — you have no idea why
				\item \textit{And no one can help you}
			\end{itemize}
		\end{block}
		\vspace{0.2cm}
		\alert{\textbf{This is normal in enterprise.}} \\
		\textit{Not exaggerating. This will be your Tuesday.}
	\end{frame}
	
	% ----------------------------------------------------------
	% SLIDE 2: THE ZERO-KNOWLEDGE DEBUGGING METHOD
	% ----------------------------------------------------------
	\begin{frame}{How to fix things when you know absolutely nothing}
		\begin{enumerate}
			\item \textbf{Stop trying to understand everything} \\
			\small{You won't. Accept it. Move to step 2.}
			
			\item \textbf{Find ONE thing that works} \\
			\small{Any old file, any past email, any working example. Copy it.}
			
			\item \textbf{Change ONE small thing} \\
			\small{If it breaks → that thing caused it. If it works → change another.}
			
			\item \textbf{Copy-paste the error into Google/StackOverflow} \\
			\small{Someone else already solved this. Steal their fix.}
			
			\item \textbf{Try a different tool} \\
			\small{TeXstudio broken? Try Overleaf. MiKTeX fails? Try TinyTeX.}
			
			\item \textbf{Reboot. Try again.} \\
			\small{You'd be surprised how often this works.}
		\end{enumerate}
		\vspace{0.2cm}
		\alert{You don't need to know WHY it works. Just make it WORK.}
	\end{frame}
	
	% ----------------------------------------------------------
	% SLIDE 3: REAL EXAMPLE — LATEX WON'T COMPILE (NO CLUE WHY)
	% ----------------------------------------------------------
	\begin{frame}{Real example: LaTeX compilation fails — zero clues}
		\begin{block}{What you see:}
			\small{\texttt{! LaTeX Error: File `somepackage.sty' not found.}} \\
			\vspace{0.1cm}
			\small{\texttt{TeXstudio: Process exited with error(s).}} \\
			\vspace{0.1cm}
			\small{\textit{You have no idea what `somepackage' is.}}
		\end{block}
		\vspace{0.2cm}
		\textbf{Zero-knowledge fix workflow:}
		\begin{enumerate}
			\item Google: \texttt{LaTeX File `xxx.sty' not found}
			\item Find answer: \texttt{\textbackslash usepackage\{xxx\}} missing or package not installed
			\item Try: Comment out that line → compiles? Yes → problem found
			\item Fix: Install package via MiKTeX console (if proxy allows)
			\item If proxy blocks it → download .sty file manually, put in same folder
			\item \textbf{Document what you did} (because future you will forget)
		\end{enumerate}
		\vspace{0.1cm}
		\alert{You didn't learn LaTeX. You learned how to survive LaTeX.}
	\end{frame}
	
	
	% ==========================================================
	% SECTION: WORKING WHEN YOU KNOW ABSOLUTELY NOTHING
	% ==========================================================
	\section{Working When You Know Nothing}
	
	% ----------------------------------------------------------
	% SLIDE 1: THE ASSIGNMENT — YOU KNOW NOTHING
	% ----------------------------------------------------------
	\begin{frame}{4.1 The assignment arrives — and you have zero information}
		\begin{block}{What you get (credit card fraud model example)}
			\begin{itemize}
				\item A ticket or email: \textit{``Update the fraud detection model documentation''}
				\item No model name — just ``fraud model''
				\item No code location — \textit{``It's in the repo''} — which repo?
				\item No database name — \textit{``Use the transaction table''} — which one? There are 40.
				\item No LaTeX path — \textit{``Docs are in the wiki''} — wiki has 3,000 pages
				\item No model website access — you don't even know if one exists
				\item No GitHub — \textit{``Ask the dev team''} — who are they?
			\end{itemize}
		\end{block}
		\vspace{0.1cm}
		\alert{\textbf{This is not an exaggeration.}} \\
		\textit{This will be a real Tuesday for you.}
	\end{frame}
	
	% ----------------------------------------------------------
	% SLIDE 2: WHAT YOU DON'T HAVE (THE FULL LIST)
	% ----------------------------------------------------------
	\begin{frame}{4.2 What you don't know — the complete inventory of missing info}
		\begin{columns}[T]
			\column{0.5\textwidth}
			\textbf{Code \& repo:}
			\begin{itemize}
				\item Where is the code? Git? Bitbucket? SVN?
				\item Repo name? No one remembers
				\item Access permissions? You have none
				\item Last commit date? Unknown
			\end{itemize}
			\vspace{0.2cm}
			\textbf{Database \& data:}
			\begin{itemize}
				\item Which database server?
				\item Table names? Unknown
				\item Schema? No documentation
				\item Data dictionary? Doesn't exist
			\end{itemize}
			\column{0.5\textwidth}
			\textbf{Documentation \& LaTeX:}
			\begin{itemize}
				\item Where are the .tex files?
				\item Wiki page? Outdated or missing
				\item SharePoint? 10 versions, no labels
				\item Model website? You can't access it
			\end{itemize}
			\vspace{0.2cm}
			\textbf{People \& process:}
			\begin{itemize}
				\item Who owns the model? Unknown
				\item Who wrote the code? Left 2 years ago
				\item Who approves changes? No idea
			\end{itemize}
		\end{columns}
		\vspace{0.2cm}
		\alert{You have nothing. But the work is still expected.}
	\end{frame}
	
	% ----------------------------------------------------------
	% SLIDE 3: HOW TO WORK WHEN YOU KNOW NOTHING (THE METHOD)
	% ----------------------------------------------------------
	\begin{frame}{4.3 The zero-knowledge workflow — start anywhere, find one thread}
		\begin{enumerate}
			\item \textbf{Find ONE person who might know something} \\
			\small{Ask: ``Who worked on fraud models last year?'' Chase that name.}
			
			\item \textbf{Find ONE file, ANY file} \\
			\small{Search shared drive for ``fraud'' + ``model''. Open everything.}
			
			\item \textbf{Find ONE database table} \\
			\small{Run: \texttt{SHOW TABLES LIKE '\%fraud\%'}. Guess. Test.}
			
			\item \textbf{Find ONE LaTeX file} \\
			\small{Search wiki for ``.tex'' or ``documentation''. Download everything.}
			
			\item \textbf{Ask ONE specific question} \\
			\small{Not ``how do I do this?'' — but ``I found file X, is that the right one?''}
			
			\item \textbf{Document what you find} \\
			\small{Create your own map: code here, data there, docs somewhere else.}
		\end{enumerate}
		\vspace{0.2cm}
		\alert{You don't need everything to start. You need ONE thread to pull.} \\
		\textit{Pull it. Something will unravel. Then you pull the next one.}
		
		\vspace{0.1cm}
		\small{This is how experienced people work when they know nothing. \\
			They don't wait for perfect info. They go find one piece and build from there.}
	\end{frame}
	
	
	% ==========================================================
	% SECTION: CROSS-TEAM CHAOS (WITHOUT ASKING ANYONE)
	% ==========================================================
% ==========================================================
% SECTION: CROSS-TEAM CHAOS (WITHOUT ASKING ANYONE)
% ==========================================================
\section{Cross-Team Chaos (No Questions Allowed)}

% ----------------------------------------------------------
% SLIDE 1: THE TEAMS — ALL SILENT, ALL UNKNOWN
% ----------------------------------------------------------
\begin{frame}{5.1 The teams that exist — but you'll never talk to them}
	\begin{block}{The official 3 lines of defense (what banks use)}
		\scriptsize
		\begin{enumerate}
			\item \textbf{Line 1: Model Development Team} — builds model, writes code, owns logic
			\item \textbf{Line 2: Model Risk / Validation (THIS IS NOT YOU)} — independent review, finds issues
			\item \textbf{Line 3: Internal Audit} — checks if everyone followed rules
		\end{enumerate}
	\end{block}
	
	\vspace{0.05cm}
	
	\begin{block}{The real teams you'll never meet}
		\scriptsize
		\begin{itemize}
			\item \textbf{Model Implementation Team} — takes code from dev, puts into production
			\item \textbf{Model Maintenance Team} — fixes bugs, updates code, deploys changes
			\item \textbf{Indian Team} — works overnight, you'll never know their names
			\item \textbf{USA Team} — works your hours, also doesn't answer emails
			\item \textbf{Database Team} — owns tables, manages access, controls schemas
			\item \textbf{Business Team} — defines requirements, uses outputs, makes decisions
		\end{itemize}
	\end{block}
	
	\vspace{0.05cm}
	\tiny\alert{Eight teams. Zero answers. You're alone.}
\end{frame}

% ----------------------------------------------------------
% SLIDE 2: WHO ARE THESE PEOPLE? (YOU'LL NEVER KNOW)
% ----------------------------------------------------------
\begin{frame}{5.2 The unknown teams — names you'll never learn}
	\begin{columns}[T]
		\column{0.5\textwidth}
		\scriptsize
		\textbf{Technical \& Implementation:}
		\begin{itemize}
			\item \textbf{Data Team} — tables, schemas, quality
			\item \textbf{Database Team} — servers, permissions
			\item \textbf{Model Dev Team} — code, logic, build
			\item \textbf{Model Implementation Team} — production
			\item \textbf{Model Maintenance Team} — patches, fixes
		\end{itemize}
		\vspace{0.1cm}
		\textbf{Offshore / Onshore:}
		\begin{itemize}
			\item \textbf{Indian Team} — night shift
			\item \textbf{USA Team} — day shift
		\end{itemize}
		
		\column{0.5\textwidth}
		\scriptsize
		\textbf{Oversight \& Business:}
		\begin{itemize}
			\item \textbf{Validation Team} — previous reviews
			\item \textbf{Audit Team} — compliance, flags
			\item \textbf{Business Team} — what model does
			\item \textbf{Product Owners} — why built
			\item \textbf{Risk Committee} — who approved
		\end{itemize}
	\end{columns}
	
	\vspace{0.1cm}
	\tiny\alert{Eleven teams. Zero names. Zero answers. You're alone.}
\end{frame}

% ----------------------------------------------------------
% SLIDE 3: HOW TO WORK WHEN YOU CAN'T ASK ANYONE
% ----------------------------------------------------------
\begin{frame}{5.3 Working alone — the no-questions-allowed method}
	\begin{block}{You cannot ask anyone. So what do you do?}
		\scriptsize
		\begin{enumerate}
			\item \textbf{Assume everything is incomplete} — docs wrong, code bugs, data gaps
			\item \textbf{Find the paper trail} — old Jira, emails, validation reports
			\item \textbf{Compare everything} — docs vs code, India vs USA handoffs
			\item \textbf{Reverse-engineer business need} — why model exists, who uses it
			\item \textbf{Make best guess, then test} — run small test, see what breaks
			\item \textbf{Document assumptions} — ``I assumed X because I found Y in old email''
		\end{enumerate}
	\end{block}
	
	\vspace{0.05cm}
	\tiny\alert{You don't need permission. One file, one guess, one test.}
\end{frame}
	
	
	
	% ----------------------------------------------------------


% ----------------------------------------------------------
% SLIDE: RUNNING NOTEBOOKS WITH NO IDEA — JUST RUN AND MESS
% ----------------------------------------------------------
\begin{frame}{5.4 Running notebooks when you know nothing — just try stuff}
	\begin{block}{The real workflow (no one explains this)}
		\tiny
		\begin{itemize}
			\item \textbf{Running code} — press play, see if it crashes
			\item \textbf{Running unit tests} — half will fail, you won't know why
			\item \textbf{Running notebook} — takes 45 minutes, then fails at the end
		\end{itemize}
	\end{block}
	
	\vspace{0.03cm}
	
	\begin{block}{The permission nightmare}
		\tiny
		\begin{itemize}
			\item \textbf{Power Broker admin account} — you don't have it, someone forgot to give it
			\item \textbf{Phone OTP permissions} — every 5 minutes, code expires, you're locked out
			\item \textbf{Driver machines} — git clone the repo, but SSH keys don't work
			\item \textbf{Slave machines} — data lives there, you can't access it
		\end{itemize}
	\end{block}
	
	\vspace{0.03cm}
	
	\begin{block}{The "just try it" method}
		\tiny
		\begin{enumerate}
			\item \textbf{Change one line} — you don't know what it does, change it anyway
			\item \textbf{Run it} — see if something different happens
			\item \textbf{If it works} — you still don't know why, but move on
			\item \textbf{If it breaks} — change it back, try something else
			\item \textbf{Repeat 47 times} — eventually it might work
		\end{enumerate}
	\end{block}
	
	\vspace{0.03cm}
	\tiny\alert{You don't need to understand. You need to try, fail, try again.}
\end{frame}


% ----------------------------------------------------------
% SLIDE: AI / COPILOT — 80% OF CODE AND DOCS WILL BE AI
% ----------------------------------------------------------
\begin{frame}{7.1 AI and Copilot — the inevitable takeover}
	\begin{block}{The reality (no one is telling you)}
		\tiny
		\begin{itemize}
			\item \textbf{Code versioning} — AI will suggest commits and diffs
			\item \textbf{Code writing} — 80\% will be generated, not typed
			\item \textbf{Unit testing} — AI writes the tests, you pray they pass
			\item \textbf{Documentation} — AI converts emails to wiki pages
			\item \textbf{LaTeX code} — AI writes the .tex, you fix the errors
		\end{itemize}
	\end{block}
	
	\vspace{0.03cm}
	
	\begin{block}{When AI helps (the good parts)}
		\tiny
		\begin{itemize}
			\item \textbf{Summarizing docs} — 50 pages → 5 bullet points
			\item \textbf{Code suggestions} — fills in the boring parts
			\item \textbf{Explaining errors} — turns garbage logs into English
			\item \textbf{First draft of anything} — you edit, not create
		\end{itemize}
	\end{block}
	
	\vspace{0.03cm}
	
	\begin{block}{When AI is dangerous (very dangerous)}
		\tiny
		\begin{itemize}
			\item \textbf{Hallucinated logic} — looks correct, is completely wrong
			\item \textbf{Data privacy issues} — you paste client data, it's gone forever
			\item \textbf{Clicks you didn't authorize} — AI runs commands you don't understand
			\item \textbf{Confident wrong answers} — AI never says "I don't know"
		\end{itemize}
	\end{block}
	
	\vspace{0.03cm}
	
	\begin{block}{Safe usage framework — "Trust but verify"}
		\tiny
		\begin{enumerate}
			\item \textbf{Use AI for first draft} — code, docs, tests, LaTeX
			\item \textbf{Assume 20\% is wrong} — always check the logic
			\item \textbf{Never paste real data} — anonymize or use fake data
			\item \textbf{Read every line before running} — AI will kill your database
			\item \textbf{Keep a human in the loop} — you are responsible, not the AI
		\end{enumerate}
	\end{block}
	
	\vspace{0.03cm}
	\tiny\alert{80\% of your job will be AI. The other 20\% is checking if AI is lying to you.}
\end{frame}

% ----------------------------------------------------------
% SLIDE 25: SUMMARY OF SKILLS
% ----------------------------------------------------------
\begin{frame}{The Ultimate Survival Toolkit}
	\begin{center}
		\small
		\begin{tabular}{p{4cm} p{7cm}}
			\toprule
			\textbf{Challenge} & \textbf{Key Skill} \\
			\midrule
			No training & Self-directed learning and building a personal knowledge base \\
			Bad documentation & Extract information and infer missing details from context \\
			Too many tools & Create a tool map and identify the primary source of truth \\
			Imposter syndrome & Track accomplishments and reframe uncertainty as learning \\
			Unclear ownership & Take temporary ownership and propose a responsible owner \\
			Office politics & Be helpful, communicate clearly, and close feedback loops \\
			Incomplete understanding & Take risk-aware actions and test assumptions incrementally \\
			\bottomrule
		\end{tabular}
	\end{center}
\end{frame}
	
	% ----------------------------------------------------------
	% SLIDE 26: MINDSET SHIFT — FROM PERFECTION TO PROGRESS
	% ----------------------------------------------------------
	\begin{frame}{The single most important mindset}
		\centering
		\Large{\textbf{Progress over perfection}}
		\vspace{0.4cm}
		
		In school: clarity → answer → grade \\
		In corporate: ambiguity → progress → iteration
		
		\vspace{0.5cm}
		\alert{The people who survive aren't the ones who know the most.} \\
		They are the ones who learn, adapt, and deliver despite chaos.
	\end{frame}
	
	% ----------------------------------------------------------
	% SLIDE 27: FINAL PRACTICAL TIPS (NEXT WEEK)
	% ----------------------------------------------------------
	\begin{frame}{What to do on Monday morning}
		\begin{enumerate}
			\item Identify ONE process that frustrates your team — document it better
			\item Create a personal  FAQ  for the top 5 questions you had this week
			\item Set up a 15-min   tool walkthrough  with a senior colleague
			\item Find one low-risk task you can own completely
			\item Write down 3 things you learned, 3 things still unclear (review with manager)
		\end{enumerate}
	\end{frame}
	
	% ----------------------------------------------------------
	% SLIDE 28: CLOSING — YOU'VE GOT THIS
	% ----------------------------------------------------------
	\begin{frame}{You're not alone — and you're more ready than you think}
		\begin{center}
			\Large{\textbf{Success isn't about knowing everything on day one.}}
			\vspace{0.3cm}
			
			It's about how effectively you learn, adapt, communicate, and solve problems when answers are unclear.
			
			\vspace{0.6cm}
			\alert{This course gave you the map. Now go walk through the chaos.}
			\vspace{0.5cm}
			
			\textit{Thank you — now go survive (and thrive) in your first corporate job.}
		\end{center}
	\end{frame}
	
	% ----------------------------------------------------------
	% SLIDE 29: APPENDIX / QR / CONTACT STYLE
	% ----------------------------------------------------------
	\begin{frame}{Additional resources}
		\textbf{Recommended follow-ups:}
		\begin{itemize}
			\item \textit{The First 90 Days} by Michael Watkins
			\item \textit{An Everyone Culture} (psychological safety)
			\item Internal wiki: Create your own  How-to  page
		\end{itemize}
		\vspace{0.3cm}
		\small{Remember: every senior person was once confused, overwhelmed, and clueless. The only difference? They kept showing up and figuring it out.}
	\end{frame}
	
	% ----------------------------------------------------------
	% SLIDE 30: THANK YOU
	% ----------------------------------------------------------
	\begin{frame}{}
		\centering
		\vspace{1cm}
		{\Huge Questions?}
		\vspace{0.5cm}
		
		\textit{Keep learning, keep shipping, keep asking smart questions.}
		\vspace{1cm}
		
		\small{© How to Survive Your First Corporate Job — all situations real, all advice tested.}
	\end{frame}
	
\end{document}